\documentclass[11pt]{article}
\usepackage[margin=1in]{geometry}
\usepackage[utf8]{inputenc}
\usepackage{amsmath,amssymb}
\usepackage{graphicx}
\usepackage{booktabs}
\usepackage{caption}
\usepackage[hidelinks]{hyperref}
\usepackage{url}
\usepackage{enumitem}
\usepackage[round,sort,authoryear]{natbib}

\newcommand{\dkl}{$\Delta$KL}

\title{\bfseries How business-model vocabulary diffuses across US industries:\\
\large Text-based evidence from trademark filings, with kill-condition tests of the diffusion structure}
\author{Eric Silver\thanks{Independent researcher; formerly a PhD candidate at Carnegie Mellon University. Comments to \texttt{epsilver@gmail.com}. The author's current employment is unrelated to this work, and the views expressed are the author's alone.}}
\date{June 2026 working paper}

\begin{document}
\maketitle

\begin{abstract}
\noindent Existing innovation measures track invention (patents, R\&D) or analyst opinion (expert ``most innovative'' lists). Both leave most service-sector and business-model innovation unmeasured. This paper applies the DeDeo prospective/retrospective Kullback--Leibler resonance framework to a USPTO trademark corpus of 2.32 million granted filings (1990--2024) across four canonical industries --- software (NICE 009), tech-services (042), transport (039), and advertising/retail (035) --- to construct a text-based novelty signal and a directly-measured cross-industry diffusion record. Three findings. First, the constructed signal \dkl{} is orthogonal-to-negative with patents within firms ($-0.05\sigma$, $t = -8$ on 14{,}463 firms), establishing it as a distinct novelty channel rather than a noisy patent proxy. Second, business-model vocabularies diffuse asymmetrically across industries: software and tech-services origin themes (cloud, AI, mobile-app, as-a-service, blockchain, streaming) arrive in transport and retail with multi-year lags of one to thirteen years, and the diffusion structure passes two independent kill-condition tests --- random (class, year) label shuffles destroy trajectory autocorrelation at $z = +46$ under LDA topic extraction and $z = +60.7$ under independently-fitted NMF, with seven of ten hand-curated seed phrases independently rediscovered by NMF as topics at $T = 50$. Third, the Schumpeter Mark~I (entrant-led innovation) hypothesis collapses on year-matched baseline contrast: raw 73.7\% entrant share among the 50 earliest carriers per theme becomes a $-2.7$pp \emph{excess} once the corpus's high overall debut rate (76.4\%) is netted out. The bounded contribution is a reproducible cross-industry text-based novelty measure that is distinct from patents, a directly-measured diffusion record with method-invariant structure, and a worked discipline result on the kind of base-rate contrast that distinguishes useful from inflated findings in trademark-text work. Code and per-firm-year panels are released for replication.
\end{abstract}

\section{Introduction}

Quantitative innovation research runs on three main rails: patent counts, patent-citation impact, and R\&D spending intensity \citep{Griliches1990, Trajtenberg1990, HallJaffeTrajtenberg2005}. All three are largely silent on the parts of the economy where most modern firm value is produced. Patents miss most service-sector activity, business-model innovation, and consumer-facing brand introductions. R\&D measurement is loose on what counts as research. Expert lists like BCG Most Innovative or MIT Technology Review's TR50 are non-replicable and survivor-biased. The standard apparatus measures invention reasonably; it does not measure most of what economists and managers mean when they say ``innovation.''

Trademark filings are public, dated, text-bearing, and filed contemporaneously with commercial intent. The goods/services descriptions applicants write to define protection scope are forced disclosures: the language must be wide enough to cover the intended commercial activity and narrow enough to be granted. Trademarks have been used for business-formation work \citep{Dinlersoz2018} and for distinguishing protective filings from substantive product introductions \citep{Graham2013}. The text of those filings has been mostly unused for innovation measurement. The closest methodological precedent in the empirical-text-economics literature is the text-based industry classification of \citet{HobergPhillips2016}, which extracts industry information from 10-K filings; my approach is methodologically related but operates on a different corpus (USPTO trademarks rather than 10-Ks), with a different signal (resonance over time rather than text similarity across firms), and at a different unit of analysis (the filing within its industry's vocabulary trajectory).

This paper applies the DeDeo prospective/retrospective Kullback--Leibler resonance framework \citep{Murdock2017, Barron2018} --- developed for measuring directional resonance in Darwin's reading notebooks and French Revolution parliamentary debates --- to a 2.32 million-filing USPTO trademark corpus covering 1990--2024 across four canonical industries: software (NICE 009), tech-services (042), transport (039), and advertising/retail (035). The selection of classes follows a canonical phenomenon the work aims to capture: ``Uber-for-X''-style transit, where business-model vocabulary born in software or tech-services arrives years later in physical-economy industries.

The paper makes three empirical claims and one methodological point. First, \dkl{} is orthogonal-to-negative with patents at the firm level: in years a firm patents more, its trademark vocabulary regresses toward class-typical rather than ahead. The within-firm coefficient is $-0.048\sigma$ ($t = -7.9$) on 14{,}463 firms with 174{,}569 firm-years. This establishes \dkl{} as a distinct novelty channel rather than a noisy patent proxy. Second, business-model vocabularies diffuse across industries asymmetrically: software and tech-services origin themes arrive in transport and retail with multi-year lags. The directional asymmetry is concrete (a transit table of fourteen phrases is provided), and the underlying diffusion structure passes two independent kill-condition tests at $z = +46$ under LDA topic extraction and $z = +60.7$ under independently-fitted non-negative matrix factorization. Third, adoption curves are Bass-fittable: of 118 (theme, class) cells with $R^2 \geq 0.7$, the bulk are platform-like (slow imitation, broad spread) rather than fad-like \citep{Bass1969}.

The methodological point is contained in a worked discipline result on the Schumpeter Mark~I~/~II debate. Raw entrant share among the 50 earliest carriers per theme is 73.7\% (48 of 50 themes entrant-dominated). Year-matched against the corpus's overall debut rate of 76.4\%, this collapses to a $-2.7$pp \emph{excess}: themes inherit the corpus's high debut rate without adding to it. The Schumpeter~I claim does not survive. This kind of base-rate contrast is the discipline I argue the trademark-text literature needs.

The bounded contribution is a reproducible cross-industry text-based novelty measure, a directly-measured cross-industry diffusion record with method-invariant structure, and a worked example of how base-rate contrast distinguishes useful from inflated findings. The paper does not claim that \dkl{} replaces patents as the canonical innovation measure (the within-firm patent orthogonality contradicts this) or that trademark text identifies a multi-year return premium (an earlier $+28$pp four-year debut excess return claim is a small-sample artifact, dropped in Section~\ref{sec:scope}); a companion paper on firm-level outcomes finds a modest, transient, R\&D-substituting margin association that I summarize but do not develop here.

Section~\ref{sec:construct} describes the construct and presents the patent-orthogonality validation. Section~\ref{sec:transit} builds the cross-class phrase-level transit table. Section~\ref{sec:kill} reports the kill-condition tests. Section~\ref{sec:bass} reports Bass fits and the software-as-net-exporter quantification. Section~\ref{sec:rq4} reports the Schumpeter~I baseline collapse. Section~\ref{sec:scope} discusses scope, limits, and the forward gesture. Section~\ref{sec:repro} is reproducibility.

\section{$\Delta$KL resonance on a USPTO trademark corpus}\label{sec:construct}

\subsection{The DeDeo resonance framework}

For a time-ordered sequence of documents, with each document at year $t$ producing a token distribution $P_t$, the DeDeo framework \citep{Murdock2017, Barron2018} defines two quantities against a window $W$ of past and future documents:

\begin{itemize}[noitemsep]
  \item \textit{Prospective novelty}: $\mathrm{KL}(P_t \,\|\, Q_{\mathrm{past}})$, where $Q_{\mathrm{past}}$ is the aggregate token distribution of documents in years $t-W$ to $t-1$.
  \item \textit{Retrospective novelty}: $\mathrm{KL}(P_t \,\|\, Q_{\mathrm{future}})$, where $Q_{\mathrm{future}}$ is the aggregate distribution of documents in years $t+1$ to $t+W$.
\end{itemize}

\noindent The signed resonance is
\[
\Delta\mathrm{KL}_t \;=\; \mathrm{KL}(P_t \,\|\, Q_{\mathrm{past}}) - \mathrm{KL}(P_t \,\|\, Q_{\mathrm{future}}).
\]

A positive $\Delta\mathrm{KL}$ marks a document whose vocabulary is unusual relative to recent past and close to near future: the field ``moved toward'' it. A negative $\Delta\mathrm{KL}$ marks the inverse, language that looked normal at filing and looked played-out in hindsight.

The framework was developed on cognitive corpora where the document author is the agent whose novelty is being measured (Darwin's reading notebooks, French Revolution speeches). A USPTO trademark goods/services description differs in that the text is typically drafted by counsel to define protection scope, not to express a state of cognition. Whether the DeDeo apparatus applied to commercially-incentivised text measures the intended construct of directional resonance is the empirical question Sections~\ref{sec:transit}--\ref{sec:rq4} take up.

\subsection{Data}

The USPTO Trademark Annual Applications backfile (product TRTYRAP) provides a complete record of US trademark filings from 1884 onward, with filing date, registration date, goods/services description, NICE class assignments, owner name, and status. For this paper I work with four NICE classes selected to exercise canonical cross-industry diffusion: software (009), scientific and technical services (042), transport (039), and advertising and retail (035). The selection follows the framing that ``Uber-for-X''-style transit --- vocabulary born in software/tech and arriving in physical-economy classes --- is the canonical phenomenon worth measuring directly.

Filing counts in each class, restricted to filings with registration\_date populated (the ``granted'' conditioning rule below), 1990 onward: software 912{,}408; tech-services 575{,}041; transport 100{,}885; advertising/retail 753{,}050. Pooled across the four classes, 2{,}317{,}899 granted filings from 1990 through 2024 enter the diffusion analysis.

\subsection{Grant conditioning}

The corpus is restricted to filings with a populated registration date. The conditioning has two motivations. First, applications include speculative or defensive filings that may not reflect realized commercial activity; the registration outcome screens out applications that did not survive examination. Second, the resonance signal is most interpretable for filings whose vocabulary reflects intended commercial use rather than maximal claim-scope drafting. Approximately 56\% of Class~009 filings in the full backfile are granted; the four-class pooled share is similar.

\subsection{Computation}

For each granted filing at year $t$ in class $c$ with goods/services text producing token distribution $P$ (unigram + bigram, after a standard analyzer that drops stop words and tokenizes within punctuation chunks), the prospective and retrospective reference distributions $Q_{\mathrm{past}}$ and $Q_{\mathrm{future}}$ are computed from all granted filings in class $c$ in the respective windows. I use $W = 5$ years (flat window), with Dirichlet smoothing ($\alpha = 1$ per cell) to avoid infinite surprise from zero-count tokens. The vocabulary is the set of unigrams and bigrams with within-class document frequency at least 50.

The construct's known fragilities have been documented in companion analysis on a complete Class~009 corpus of 1{,}807{,}636 filings. The prospective and retrospective KL are correlated at $r = +0.971$ across 1.56M scored filings; signed \dkl{} is therefore a small residual of two large, highly correlated quantities and is noisy at the per-filing level. The magnitude $|\Delta\mathrm{KL}|$ inflates monotonically for short filings: $0.293$ for documents under 5 terms, decaying to $0.147$ for documents of 40 or more terms. Pre-1995 readings are contaminated by thin past (mean \dkl{} of $+2.05$ in 1990). Window sensitivity is bounded: across $W \in \{3, 5, 7\}$ years, only 6.9--10.4\% of filings flip sign --- the direction of the residual is robust to half-life kernel choice, the magnitude is kernel-relative.

The aggregation level at which the construct is interpretable is taken up below. The diffusion analyses in Sections~\ref{sec:transit}--\ref{sec:rq4} use theme-level aggregation; the construct-distinctness analysis in the next subsection uses firm-year aggregation.

\subsection{$\Delta$KL is orthogonal-to-negative with patents within firm}\label{sec:patent_check}

If trademark \dkl{} and patent counts were noisy measures of the same underlying invention construct, they should correlate within firms over time. They do not.

The within-firm analysis uses a published panel of 174{,}569 firm-years covering 14{,}463 trademark owners that have been name-matched to PatentsView patent assignees \citep{PatentsView}. The panel is constructed in companion work and is used here only to validate the construct against the patent-track measure. I fit a two-way (firm + year) fixed-effects regression of standardized firm-year mean \dkl{} on $\log(1 + \mathrm{patents})$ with firm-clustered standard errors and $\sigma(\Delta\mathrm{KL}) = 0.214$ on the collapsed panel.

\begin{table}[h]\centering\small
\caption{Within-firm \dkl{} versus patenting. Coefficients in standard deviations of \dkl{} per unit $\log(1 + \mathrm{patents})$.}
\label{tab:patent_within_firm}
\begin{tabular}{lrrr}
\toprule
Specification & coef ($\sigma$) & SE & $t$ \\
\midrule
Between-firm correlation (firm means) & $-0.020$ & --- & --- \\
Within-firm correlation (firm+year demeaned) & $-0.015$ & --- & --- \\
Two-way FE, contemporaneous & $-0.048$ & $0.006$ & $-7.9$ \\
\quad robustness: patents = max over matched assignees & $-0.048$ & $0.006$ & $-8.0$ \\
\quad controlling $\log(\mathrm{n\_filings})$ & $-0.051$ & $0.006$ & $-8.4$ \\
Lead/lag: patents$_{t-1}$ & $-0.054$ & $0.009$ & $-5.8$ \\
\quad\quad\quad\ \ patents$_{t}$ & $-0.019$ & $0.010$ & $-1.9$ \\
\quad\quad\quad\ \ patents$_{t+1}$ & $+0.010$ & $0.010$ & $+1.0$ \\
\bottomrule
\end{tabular}
\end{table}

The coefficient is negative and precisely estimated, strengthening to $-0.051$ under filing-count control. The lag structure is informative: the most negative coefficient sits on patents$_{t-1}$. In years a firm patents more, the firm's trademark vocabulary regresses toward class-typical the following year rather than ahead. Cross-sectionally, the Spearman rank correlation between \dkl{} and patent count is $-0.015$ ($p = 8 \times 10^{-10}$) on the 174{,}569 firm-years.

The interpretation: \dkl{} is not measuring patent-track invention. Within firms across time, the two measures move in opposite directions rather than orthogonally. The construct earns the claim that it measures a distinct channel of novelty rather than a noisy patent proxy. This patent-orthogonality result is also useful as protection against an alternative reading of subsequent results: if the diffusion structure in Sections~\ref{sec:transit}--\ref{sec:rq4} reflected something that patents also measured, the present result would be inconsistent with that.

\subsection{Discriminant validation against expert ``most innovative'' lists}

A brief out-of-corpus validation against expert ``most innovative'' rankings is informative but small-$n$. Firms appearing on BCG Most Innovative (2021--23) score $+0.29$ population-$\sigma$ above the non-listed baseline (Mann--Whitney $p = 0.012$, $n = 31$ matched to a CIK in our public panel). MIT Technology Review TR50 (2015) firms score $+0.50\sigma$ above ($p = 0.0006$, $n = 17$). Pooled (union of both lists): $+0.34\sigma$, $p = 0.0003$, $n = 43$. The lists were built without using trademark text.

This is the first out-of-corpus validation in the project. The sample is genuinely small, the lists are survivor-biased (current firms), and the temporal alignment is loose (each firm's all-years \dkl{} mean against a recent list). I would not stake a strong construct claim on $n = 48$ firms. The result is consistent with \dkl{} discriminating expert-perceived market-facing innovation while being orthogonal to patent-based technical invention --- two different innovation constructs rather than the same one measured imperfectly. This complements the within-firm patent-orthogonality result above.

\section{Cross-industry theme transit at the phrase level}\label{sec:transit}

\subsection{Operationalization}

The proposal frames the canonical phenomenon as transit of business-model vocabulary across industries: an ``Uber-for-X'' description born in software/tech and arriving years later in transport, retail, or another physical-economy class. To make this operational without committing to a specific topic-extraction method --- which the kill-condition tests in Section~\ref{sec:kill} take up --- I define a list of 14 multi-token business-model phrases and track their prevalence across (class, year) cells directly.

For each phrase, I compute per (class, year) the share of granted filings whose goods/services text contains the phrase via case-insensitive regex match on the raw text. This bypasses the analyzer's stopword removal, which would mangle multi-word phrases such as ``on demand'' or ``as a service.'' I then record the first year in which each class crosses a 1\% prevalence floor as the ``arrival'' year for that theme in that class.

\subsection{The transit table}

\begin{table}[h]\centering\small
\caption{First year a theme reaches 1\% prevalence among granted filings, by class. ``--'' = never clears the floor.}
\label{tab:transit}
\begin{tabular}{lrrrr}
\toprule
Theme & software & tech-svcs & transport & advert/retail \\
\midrule
as-a-service           & 2011 & \textbf{2009} & 2013 & 2012 \\
mobile-app             & 2010 & \textbf{2009} & 2012 & 2012 \\
cloud                  & 2012 & \textbf{2010} & 2011 & 2014 \\
streaming              & \textbf{2008} & 2008 & 2013 & 2014 \\
artificial-intelligence & 2017 & \textbf{2016} & 2018 & 2018 \\
augmented/virtual-reality & \textbf{2015} & 2016 & 2022 & 2022 \\
blockchain / crypto    & 2021 & \textbf{2018} & --   & 2022 \\
real-time              & \textbf{2001} & 2003 & 2009 & 2014 \\
on-demand              & 2017 & 2016 & \textbf{2014} & 2017 \\
\bottomrule
\end{tabular}
\end{table}

Of the phrases that clear the 1\% floor in transport and advertising/retail, almost all originate in software (009) or tech-services (042) and arrive in the physical-economy classes with lags of one to thirteen years. The exception is on-demand, which originates in transport (2014) before software/tech (2016--17) --- a genuine counter-example worth flagging. Some phrases (peer-to-peer, IoT, marketplace) do not clear the 1\% floor at exact-phrase regex in some classes; the directional pattern below is robust to including or excluding these.

\subsection{Caveats}

The 1\% floor is sensitive to class size. Transport (039) has roughly one-eighth the filing volume of software (009), so the floor trips on far fewer absolute filings. A phrase below 1\% in transport may still be present with substantive absolute volume. The 1\% choice is conventional; a different floor (0.5\% or 2\%) shifts arrival years by a year or two but does not reverse the directional pattern.

The asymmetry is also a four-class result. A full 45-class build with NICE-to-NAICS crosswalk would test whether software-as-net-exporter is a class-size artifact or a substantive directional pattern. The current evidence is consistent with the directional reading but cannot exclude a corpus-construction alternative.

What the result earns: on a substantively large four-class corpus with grant conditioning, business-model phrases originate in software/tech-services and arrive years later in transport and retail more often than the reverse. This claim is replicable from the published code.

\section{Kill conditions: the diffusion structure is method-invariant}\label{sec:kill}

The transit table in Section~\ref{sec:transit} depends on the curated phrase list. A different list, or a fully data-driven topic extraction, could produce a different picture. This section tests whether the diffusion structure is a property of the corpus or an artifact of either the phrase choice or the topic-extraction method.

\subsection{LDA topic extraction}

I fit a Latent Dirichlet Allocation topic model with $T = 50$ topics on a stratified sample of 200{,}000 granted filings (50{,}000 per class, capped at class size; transport is smaller after sampling). The vocabulary is unigrams + bigrams with document frequency $\geq 200$, giving $V = 69{,}361$ terms. The fitted model is transformed across the full 2.32 million-filing pool to assign each filing a topic distribution and a top topic.

The resulting top-word lists are recognizable business-model components. Selected examples: T0 \{health, care, wellness, healthcare\}; T1 \{autonomous, loading, unloading, truck, material handling\}; T2 \{store, retail, wholesale\}; T10 \{electronic transmission, computer networks, data, communication\}; T30 \{online, website, non-downloadable, information\}; T49 \{vehicles, traffic, bicycle, lights\}. I make no claim that 50 unsupervised topics constitute the final theme set --- cross-method stability and human face-validation remain owed before any individual theme is treated as headline. The kill-condition tests below are the cheap go/no-go check on whether \emph{any} theme extraction can support the diffusion claim.

\subsection{Shuffle null on theme trajectories}

For each (theme, class) prevalence trajectory over 1990--2024 --- where prevalence is the share of granted filings in that class-year whose top topic is that theme --- I compute the lag-1 temporal autocorrelation. Real themes should ramp up smoothly within classes (positive autocorrelation); random (class, year) label permutations should average near zero. I randomly permute each filing's (class, year) labels four times and recompute the autocorrelation.

\begin{table}[h]\centering\small
\caption{Lag-1 autocorrelation of (theme, class) prevalence trajectories under LDA top-topic assignment.}
\label{tab:lda_null}
\begin{tabular}{lcc}
\toprule
& mean autocorrelation & (theme, class) cells \\
\midrule
Real panel & $+0.687$ & 200 \\
Shuffle 1 & $-0.041$ & 200 \\
Shuffle 2 & $-0.041$ & 200 \\
Shuffle 3 & $-0.003$ & 200 \\
Shuffle 4 & $-0.024$ & 200 \\
\midrule
Shuffle mean $\pm$ sd & $-0.027 \pm 0.016$ & --- \\
\multicolumn{2}{l}{$z = (\mathrm{real} - \mathrm{shuffle})/\mathrm{sd}$} & $+46$ \\
\bottomrule
\end{tabular}
\end{table}

The real panel's themes ramp up smoothly within their classes; the shuffled panel is noise around zero. The null is rejected at any plausible significance level. The diffusion structure is not an artifact of arbitrary panel slicing.

\subsection{Independent NMF replication}

LDA is a generative topic model with specific distributional assumptions. To test whether the diffusion structure depends on these assumptions, I refit topic structure with non-negative matrix factorization (NMF) --- a different method, different objective, different convergence path --- on the same 200{,}000-filing sample and vocabulary, with $T = 50$, then transform across the full pool.

The two models partition the corpus differently in detail. The adjusted Rand index between LDA top-topic and NMF top-topic on a 500{,}000-filing sample is $+0.214$: modest but well above random partition. Greedy Jaccard matching on top-10 words gives 28\% of NMF topics aligning with an LDA topic at Jaccard $\geq 0.3$ and 14\% at $\geq 0.5$. The strongly-stable intersection includes textbook business-model components (energy/renewable; apparel; retail/store services; entertainment/music; control systems/sensors) at Jaccard $\geq 0.67$. The 72--86\% of LDA topics without a strongly-aligned NMF counterpart should be treated as method-dependent and not promoted to headline.

The kill condition on the NMF panel:

\begin{table}[h]\centering\small
\caption{Shuffle null replicated on NMF top-topic panel.}
\label{tab:nmf_null}
\begin{tabular}{lc}
\toprule
& mean autocorrelation \\
\midrule
Real panel (NMF) & $+0.689$ \\
Shuffle mean $\pm$ sd & $-0.021 \pm 0.012$ \\
$z$ & $+60.7$ \\
\bottomrule
\end{tabular}
\end{table}

The NMF real autocorrelation is $+0.689$ --- essentially identical to LDA's $+0.687$ --- and the shuffle baseline is $-0.021 \pm 0.012$, giving $z = +60.7$, more decisive than the LDA result. The diffusion structure is method-invariant: the trajectories are properties of the corpus, not artifacts of either topic model.

\subsection{Seed-phrase concordance under NMF}

A direct check on whether NMF independently rediscovers the curated phrases of Section~\ref{sec:transit}: for each phrase, I identify the NMF topic whose top-50 words most contain the phrase variants.

\begin{table}[h]\centering\small
\caption{Seed phrase $\rightarrow$ best-matching NMF topic.}
\label{tab:seed_concordance}
\begin{tabular}{ll}
\toprule
seed phrase & NMF topic (top-50 words contain) \\
\midrule
blockchain    & NMF[19]: cryptocurrency, blockchain \\
cloud         & NMF[2]: cloud computing \\
AI / ML       & NMF[16]: artificial intelligence, machine learning \\
SaaS / `as-a-service' & NMF[12]: software service \\
mobile-app    & NMF[48]: mobile applications \\
AR / VR       & NMF[45]: virtual reality, augmented reality \\
marketplace   & NMF[45]: blended with AR/VR \\
streaming     & --- no clean NMF topic \\
IoT           & --- no clean NMF topic \\
on-demand     & --- no clean NMF topic \\
\bottomrule
\end{tabular}
\end{table}

Seven of the ten seed phrases have a clear NMF-topic correspondence. The three missing (streaming, IoT, on-demand) are either too generic to concentrate in a single topic at $T = 50$ or too brief at the phrase level to provide a discriminating bag-of-words signal. Increasing $T$ or using phrase-emphasized vocabulary could capture these; the absence is not damaging to the kill-condition argument.

\subsection{What the kill conditions earn}

The kill conditions establish that the cross-class diffusion structure is real in the corpus, not an artifact of phrase choice, label assignment, or topic extraction. They do not establish that any individual LDA or NMF topic is a substantively meaningful business-model theme; the cross-method stability filter excludes 72--86\% of LDA topics from candidate-headline status. The diffusion structure stands; the headline theme list does not, until human face-validation is added in follow-up work.

\section{Adoption shape and the software-as-net-exporter quantification}\label{sec:bass}

\subsection{Bass model fits}

For each (theme, class) cell where the theme reaches the 1\% prevalence floor and is tracked for at least six years post-arrival, I fit the Bass (1969) diffusion model to the normalized cumulative prevalence. The Bass model parameterizes adoption with an innovation coefficient $p$ (rate at which non-adopters become adopters independent of social pressure), an imitation coefficient $q$ (rate at which adoption increases with the cumulative adopter share), and a saturation level $m$. I fit $p$ and $q$ on the normalized cumulative share with non-linear least squares, retaining fits with $R^2 \geq 0.7$.

\begin{table}[h]\centering\small
\caption{Bass fits across (theme, class) cells.}
\label{tab:bass}
\begin{tabular}{lr}
\toprule
Statistic & Value \\
\midrule
Fits retained ($R^2 \geq 0.7$) & 118 \\
Median $p$ (innovation coefficient) & 0.017 \\
Median $q$ (imitation coefficient) & 0.114 \\
Median $q/p$ & $\sim 7$ \\
``Platform-like'' ($q/p < 20$) & 88 \\
``Fad-like'' ($q/p \geq 50$) & 9 \\
\bottomrule
\end{tabular}
\end{table}

The median $q/p \approx 7$ falls in the range typically reported for consumer-product Bass fits in the diffusion-of-innovations literature, lending external plausibility to the (theme, class) adoption curves as real adoption events. The bulk of the fits are platform-like --- broad gradual adoption with moderate word-of-mouth --- and a small tail of 9 of 118 ($\sim 8\%$) shows fad-shaped curves with high $q/p$ indicating sharp imitation pickup. That the curves are fittable at all is consistent with the diffusion structure of Section~\ref{sec:kill}.

\subsection{Software/tech as net exporter, quantified}

The phrase-level transit in Section~\ref{sec:transit} showed software/tech-as-exporter qualitatively. Using the 50 LDA themes, I count themes that reach at least the 1\% floor in at least two classes and identify the origin (earliest class to cross the floor) and the second class (next class to cross). This produces a directed graph of theme flow between the four classes.

\begin{table}[h]\centering\small
\caption{Theme flow origin $\rightarrow$ second-arrival class.}
\label{tab:flow}
\begin{tabular}{lr}
\toprule
Origin $\rightarrow$ second class & themes \\
\midrule
software $\rightarrow$ tech-svcs & 14 \\
software $\rightarrow$ transport & 11 \\
software $\rightarrow$ advert/retail & 8 \\
transport $\rightarrow$ tech-svcs & 7 \\
transport $\rightarrow$ advert/retail & 6 \\
transport $\rightarrow$ software & 5 \\
tech-svcs $\rightarrow$ software & 4 \\
advert/retail $\rightarrow$ tech-svcs & 5 \\
\bottomrule
\end{tabular}
\end{table}

Of 33 themes reaching at least two classes, 33 outflow edges originate in software (009 and 042 combined) and only 4 return there. Software and tech-services are the net exporters of business-model vocabulary to the other two classes; the return flow is small. This confirms the phrase-level asymmetry at the theme level.

With only four nodes a centrality estimate of the kind discussed in the product-space literature \citep{HidalgoHausmann2009} is premature. A 45-class build would let me locate each class in a business-model space and ask whether centrality predicts theme-birth rate, theme-export volume, or theme-import variety. That extension is the natural follow-up paper.

\section{A worked discipline result: Schumpeter Mark~I collapses on baseline contrast}\label{sec:rq4}

The classical Schumpeter Mark~I~/~II debate asks whether new innovations are predominantly born from new entrants (Mark~I) or from incumbent diversification (Mark~II) \citep{Schumpeter1942}. The standard empirical form is to identify a population of early adopters of an innovation, classify them as entrant or incumbent, and report the entrant share.

I apply this standard form to the 50 LDA themes. For each theme, I take its 50 earliest carriers (sorted by filing year, conditional on top-topic weight $\geq 0.20$). I classify each carrier's owner as an entrant if its first granted mark across the four-class corpus is in the carrier-filing year or one year earlier --- that is, the carrier is filing at or near the owner's debut.

\begin{table}[h]\centering\small
\caption{Raw entrant share among the 50 earliest carriers per theme.}
\label{tab:raw_schumpeter}
\begin{tabular}{lr}
\toprule
Statistic & Value \\
\midrule
Themes scored & 50 \\
Mean entrant share among 50 earliest carriers & 73.7\% \\
Median & 73.0\% \\
Themes entrant-dominated ($\geq 50\%$ entrants) & 48 / 50 \\
Themes incumbent-dominated ($\leq 20\%$ entrants) & 0 / 50 \\
\bottomrule
\end{tabular}
\end{table}

The raw numbers look like strong Schumpeter Mark~I: 48 of 50 themes are entrant-dominated and zero are incumbent-dominated. This is the result that does not survive scrutiny.

The corpus's overall debut rate among granted filings in 1990--2024 is 68--77\% per year across the four classes. With a year-weighted mean (weighting by each theme's earliest-carrier year distribution), the baseline is 76.4\%. \emph{Most filings in the four-class corpus over this period are by owners debuting that year.} The corpus is mechanically saturated in entrants.

The disciplined version of the analysis takes the difference between a theme's entrant share and the year-matched population baseline:

\begin{table}[h]\centering\small
\caption{RQ4 with year-matched population baseline.}
\label{tab:rq4_baseline}
\begin{tabular}{lr}
\toprule
Statistic & Value \\
\midrule
Theme entrant share (mean / median) & 73.7\% / 73.0\% \\
Year-matched baseline (mean / median) & 76.4\% / 76.7\% \\
\textbf{Excess (theme $-$ baseline)} & $\mathbf{-2.7}$pp / $-3.7$pp \\
Themes with positive excess & 22 / 50 \\
Themes with excess $> 10$pp & 9 / 50 \\
Themes with negative excess & 28 / 50 \\
\bottomrule
\end{tabular}
\end{table}

The mean theme excess is $-2.7$ percentage points, the median $-3.7$. Twenty-eight of 50 themes are \emph{below} baseline. Only nine themes show real entrant excess above 10pp. The Schumpeter Mark~I claim collapses: themes inherit the corpus's high debut rate without adding to it. The disciplined finding is that the typical theme is not disproportionately entrant-born once we account for who files in the relevant years.

The methodological point worth stating: this kind of base-rate contrast is what distinguishes useful from inflated findings in trademark-text work. The same flaw will appear elsewhere --- anywhere a corpus has a non-uniform base rate over the relevant population dimension, raw theme-population shares can mislead. RQ4 was the obvious instance. The discipline is to look for the analogous flaw systematically in future analyses. The 9 themes with real positive excess (in our case, several blockchain/crypto and AI-era themes in their early years) merit closer examination in follow-up work as candidate genuinely entrant-led innovations, distinct from the typical theme.

\section{Scope, limits, and a forward gesture}\label{sec:scope}

\subsection{What this paper claims and does not claim}

The paper makes the following claims and not others.

\dkl{} is a reproducible text-based novelty signal on USPTO trademark goods/services text, validated against patents (orthogonal-to-negative within firm) and against expert ``most innovative'' lists (small-$n$ positive). The construct is distinct from patent-track invention; the strong-sense ``\dkl{} measures innovation'' claim is not licensed by the within-firm patent-orthogonality result and is not made.

The theme-level diffusion structure is real in the corpus and method-invariant under LDA and NMF extraction ($z = +46$ / $+60.7$ on shuffle null). Cross-industry transit of business-model vocabulary is asymmetric, with software/tech as net exporter to physical-economy classes on multi-year lags.

The Schumpeter Mark~I result does not survive year-matched baseline contrast; the disciplined finding is that themes do not produce excess entrant share once the corpus's high overall debut rate is accounted for.

What the paper does \emph{not} claim:
\begin{itemize}[noitemsep]
\item That \dkl{} replaces patents as the canonical innovation measure. The within-firm patent-orthogonality result refutes this.
\item That trademark text identifies a persistent return premium. An earlier $+28$pp four-year debut excess return claim is a small-sample artifact --- its horizon profile is non-monotone (peaks at 4y, collapses at 5y) and the debut/firm-year ratio grows monotonically with horizon, the signature of right-tail winners in a 1{,}447-firm debut sample, not a return premium. The claim is dropped.
\item That the diffusion record causally identifies an industry's propensity to import vs.\ export business-model vocabulary. The four-class subset is consistent with the directional reading but cannot exclude class-size or corpus-construction alternatives. A 45-class build with NICE-NAICS crosswalk is the natural extension.
\item That the kill-condition tests establish individual themes as substantively meaningful. The cross-method stability filter excludes 72--86\% of LDA topics; the headline theme set requires human face-validation.
\end{itemize}

\subsection{Forward gesture: firm-level outcomes}

A companion paper takes up the firm-level outcomes question that this paper does not. Using the SEC Financial Statement Data Sets firm-year panel (2009--2024; 83{,}095 firm-years; 13{,}127 unique CIKs) and a USPTO-owner-to-CIK crosswalk (7{,}233 exact matches), within-firm two-way fixed-effects analysis finds that, per $\sigma$ of within-firm borrowed-novelty (a high-within / low-cross resonance gap), firms experience a contemporaneous $+1.4$pp gross margin lift ($t = +2.1$) that fades by $t+1$ and an offsetting cross-sectional reduction in R\&D intensity. The relationship between text-based novelty and business success is, on this evidence, modest, contemporaneous, transient, and substitutive with R\&D rather than complementary.

This is a real but bounded outcome relationship that the companion paper develops in detail. For the present paper, it functions as a forward gesture: the diffusion claims here do not depend on the firm-outcome result, and the firm-outcome result does not depend on the diffusion claims. The two papers are complementary, not nested.

\subsection{What would falsify the diffusion claims}

I would update on the diffusion claim if:
\begin{itemize}[noitemsep]
\item The cross-class transit asymmetry reversed under a 45-class build with proper class-size controls and NICE--NAICS crosswalk.
\item The kill-condition shuffle null failed to be rejected under a topic-extraction method structurally different from LDA and NMF --- e.g., a sentence-embedding-cluster approach or graph-based community detection. The current LDA and NMF are both bag-of-words matrix factorizations on the same vocabulary; the cross-paradigm test is owed.
\item The seed-phrase rediscovery under NMF dropped well below 4 of 10 at $T = 50$ with this vocabulary, suggesting NMF was not finding the same theme structure LDA does at this scale.
\end{itemize}

I would update on the patent-orthogonality claim if the within-firm coefficient reversed under value-weighting or under an alternative trademark-owner-to-patent-assignee match (an alternative to PatentsView).

\section{Reproducibility}\label{sec:repro}

Code, intermediate data panels, and the integrated working report are available at \url{https://github.com/ericsilver/novelty}. The full pipeline rebuilds from a USPTO Open Data Portal API key with a single \texttt{make tm} command for the four-class corpus build (multi-hour on commodity hardware; the underlying TRTYRAP backfile is approximately 12 GB streamed once). Analysis scripts and outputs:

\begin{itemize}[noitemsep]
\item Within-firm patent validation: \texttt{scripts/wsC\_within\_firm\_patents.py} $\rightarrow$ \texttt{paper/results/wsC\_within\_firm\_patents.json}
\item Expert-list discriminant: \texttt{scripts/wsD\_validity\_battery.py} $\rightarrow$ \texttt{paper/results/wsD\_validity\_battery.json}
\item Cross-class phrase transit: \texttt{scripts/diff2\_theme\_transit.py} $\rightarrow$ \texttt{paper/results/diff2\_transit.json}
\item LDA + kill condition: \texttt{scripts/diff3\_lda\_themes.py} $\rightarrow$ \texttt{paper/results/diff3\_themes.json}, \\
   \texttt{data/processed/\{filing\_topics, theme\_panel\}.parquet}
\item Bass fits and flow: \texttt{scripts/diff4\_phase1\_rqs.py} $\rightarrow$ \texttt{paper/results/diff4\_phase1.json}
\item RQ4 baseline contrast: \texttt{scripts/diff5\_rq4\_baseline.py} $\rightarrow$ \texttt{paper/results/diff5\_rq4\_baseline.json}
\item NMF replication: \texttt{scripts/diff6\_nmf\_compare.py} $\rightarrow$ \texttt{paper/results/diff6\_nmf.json}, \\
   \texttt{data/processed/filing\_topics\_nmf.parquet}
\end{itemize}

Companion firm-outcome analyses (borrowed novelty, demand-pull association, face-validation form) are in \texttt{diff7--9} and are summarized in the companion paper. The full analysis runs under Python 3.12 with \texttt{polars}, \texttt{numpy}, \texttt{scipy}, \texttt{statsmodels}, \texttt{scikit-learn}, and \texttt{matplotlib}.

\begin{thebibliography}{99}

\bibitem[Barron et al.(2018)]{Barron2018}
Barron, A.~T., Huang, J., Spang, R.~L., \& DeDeo, S. (2018). Individuals, institutions, and innovation in the debates of the French Revolution. \emph{Proceedings of the National Academy of Sciences}, 115(18), 4607--4612.

\bibitem[Bass(1969)]{Bass1969}
Bass, F.~M. (1969). A new product growth for model consumer durables. \emph{Management Science}, 15(5), 215--227.

\bibitem[Dinlersoz et al.(2018)]{Dinlersoz2018}
Dinlersoz, E., Goldschlag, N., Myers, A., \& Zolas, N. (2018). An anatomy of US firms seeking trademark registration. \emph{NBER Working Paper} No.\ 25038.

\bibitem[Graham et al.(2013)]{Graham2013}
Graham, S.~J.~H., Hancock, G., Marco, A.~C., \& Myers, A. (2013). The USPTO trademark case files dataset: Descriptions, lessons, and insights. \emph{Journal of Economics \& Management Strategy}, 22(4), 669--705.

\bibitem[Griliches(1990)]{Griliches1990}
Griliches, Z. (1990). Patent statistics as economic indicators: A survey. \emph{Journal of Economic Literature}, 28(4), 1661--1707.

\bibitem[Hall et al.(2005)]{HallJaffeTrajtenberg2005}
Hall, B.~H., Jaffe, A., \& Trajtenberg, M. (2005). Market value and patent citations. \emph{RAND Journal of Economics}, 36(1), 16--38.

\bibitem[Hidalgo \& Hausmann(2009)]{HidalgoHausmann2009}
Hidalgo, C.~A., \& Hausmann, R. (2009). The building blocks of economic complexity. \emph{Proceedings of the National Academy of Sciences}, 106(26), 10570--10575.

\bibitem[Hoberg \& Phillips(2016)]{HobergPhillips2016}
Hoberg, G., \& Phillips, G. (2016). Text-based network industries and endogenous product differentiation. \emph{Journal of Political Economy}, 124(5), 1423--1465.

\bibitem[Murdock et al.(2017)]{Murdock2017}
Murdock, J., Allen, C., \& DeDeo, S. (2017). Exploration and exploitation of Victorian science in Darwin's reading notebooks. \emph{Cognition}, 159, 117--126.

\bibitem[PatentsView(2024)]{PatentsView}
PatentsView. (2024). Patent and assignee data, US Patent and Trademark Office. \texttt{https://patentsview.org}.

\bibitem[Schumpeter(1942)]{Schumpeter1942}
Schumpeter, J.~A. (1942). \emph{Capitalism, Socialism and Democracy}. Harper \& Brothers.

\bibitem[Trajtenberg(1990)]{Trajtenberg1990}
Trajtenberg, M. (1990). A penny for your quotes: Patent citations and the value of innovations. \emph{RAND Journal of Economics}, 21(1), 172--187.

\end{thebibliography}

\end{document}
